\chapter{libfabric and Verbs Call Reference}
\label[appendix]{app:calls}

This appendix lists the pool's own interface and the library calls behind
it, read from the source of the two implementations \citep{scuffedrdma}. It
is a reference for \cref{ch:libfabric,ch:process}.

\section{The Interface of the Pool}

One C header declares the interface. \Cref{tab:sfapi} gives its calls. A
class is hot or cold, and offsets are positions inside the registered arena.

\begin{tightblock}
\begin{longtable}{L{4.1cm}L{10.3cm}}
  \caption[The pool's interface]{The pool's interface.}\label{tab:sfapi}\\
  \toprule
  Call & What it does \\
  \midrule
  \endfirsthead
  \toprule
  Call & What it does \\
  \midrule
  \endhead
  \code{sf_opts_default} & fills the options with defaults: provider and endpoint type per class, queue depths, arena size, completion method for the hot class, blocking wait for the cold class \\
  \code{sf_open} & opens the fabric, domain, endpoints and queues for both classes and registers the arena \\
  \code{sf_listen} & server side: waits for one peer on a TCP side channel, exchanges addresses and keys, accepts the connections; endpoints without a connection exchange one message, so that their first use is complete \\
  \code{sf_connect} & client side: connects the side channel and both endpoints \\
  \code{sf_close} & closes everything in reverse order \\
  \code{sf_buf}, \code{sf_buf_len} & start and length of the data region of the arena \\
  \code{sf_write} & one-sided write from a local offset to a remote offset on one class; flags ask for immediate data, for a copy at the call when the payload is small, or hint that more operations follow \\
  \code{sf_read} & one-sided read from a remote offset into a local offset \\
  \code{sf_send} & small two-sided message, with or without immediate data \\
  \code{sf_poll_rx} & returns arrivals on a class: immediate values of writes, and messages that landed in receive slots; posts the consumed receives again \\
  \code{sf_progress} & reads the transmit completion queue or counter of a class without blocking \\
  \code{sf_wait_tx} & waits until the number of completed operations on a class reaches a target, or a timeout \\
  \code{sf_tx_posted}, \code{sf_tx_done} & operations posted and completed on a class so far \\
  \code{sf_rma_events} & number of remote writes that have landed, where the provider counts them \\
  \code{sf_inject_size} & largest payload the provider copies at the call \\
  \code{sf_cq_data_size} & bytes of immediate data the provider carries \\
  \code{sf_describe} & text description of what was opened for a class: provider, endpoint type, device \\
  \code{sf_capable} & whether a class supports a counter, remote-write events, ordered writes, or a blocking wait \\
  \code{sf_oob_barrier}, \code{sf_oob_xchg}, \code{sf_oob_fd} & helpers for benchmarks on the side channel: a barrier, an exchange of one 64-bit value, and the socket itself \\
  \code{sf_strerror}, \code{sf_now_ns} & error text, and a clock in nanoseconds \\
  \bottomrule
\end{longtable}
\end{tightblock}

\section{Library Calls by Step}

\Cref{tab:calls} follows the phases of \cref{ch:process} and names the calls
each implementation makes.

\begin{tightblock}
\begin{longtable}{L{2.9cm}L{5.7cm}L{5.6cm}}
  \caption[Library calls made by the two implementations, by step]{Library calls made by the two implementations, by step.}\label{tab:calls}\\
  \toprule
  Step & libscuffedfabric (libfabric) & libscuffedverbs (verbs) \\
  \midrule
  \endfirsthead
  \toprule
  Step & libscuffedfabric (libfabric) & libscuffedverbs (verbs) \\
  \midrule
  \endhead
  Find and open the device & \code{fi_allocinfo}, \code{fi_getinfo}, \code{fi_fabric}, \code{fi_domain}, \code{fi_freeinfo} & \code{ibv_get_device_list}, \code{ibv_open_device}, \code{ibv_query_port}, \code{ibv_query_gid}, \code{ibv_alloc_pd} \\
  Register the arena & \code{fi_mr_reg}; \code{fi_mr_bind} and \code{fi_mr_enable} where the provider asks; \code{fi_mr_desc}, \code{fi_mr_key} & \code{ibv_reg_mr} \\
  Completion queues and counters & \code{fi_cq_open}, \code{fi_cntr_open} & \code{ibv_create_cq}, \code{ibv_create_comp_channel} \\
  Endpoints & \code{fi_endpoint}, \code{fi_ep_bind}, \code{fi_enable} & \code{ibv_create_qp} \\
  Connect, connected endpoints & \code{fi_eq_open}, \code{fi_passive_ep}, \code{fi_pep_bind}, \code{fi_listen}, \code{fi_connect}, \code{fi_accept}, \code{fi_eq_sread} & \code{ibv_modify_qp} three times, to initialized, ready to receive and ready to send, with the peer's values from the side channel \\
  Connect, datagram endpoints & \code{fi_av_open}, \code{fi_av_insert}, \code{fi_getname} & not offered \\
  Post receives & \code{fi_recv} & \code{ibv_post_recv} \\
  Write, read, send & \code{fi_writemsg}, \code{fi_readmsg}, \code{fi_sendmsg}, with flags for immediate data, a copy at the call, and batching & \code{ibv_post_send} with the opcode for a write, a write with immediate data, a read or a send \\
  Poll for completions & \code{fi_cq_read}, \code{fi_cntr_read} & \code{ibv_poll_cq} \\
  Blocking wait & \code{fi_cq_sread} & \code{ibv_req_notify_cq}, \code{ibv_get_cq_event}, \code{ibv_ack_cq_events} \\
  Errors & \code{fi_cq_readerr}, \code{fi_eq_readerr}, \code{fi_cntr_readerr}, and the matching text calls & the status field of every completion, \code{ibv_wc_status_str} \\
  Close & \code{fi_close} on every object & \code{ibv_destroy_qp}, \code{ibv_destroy_cq}, \code{ibv_destroy_comp_channel}, \code{ibv_dereg_mr}, \code{ibv_dealloc_pd}, \code{ibv_close_device} \\
  \bottomrule
\end{longtable}
\end{tightblock}

Three details of the libfabric implementation are easily overlooked. Classes that
use the same provider and endpoint type share one fabric, one domain and one
memory registration, which is the libfabric form of two queue pairs on one
protection domain. With the counter selected for the hot class, the endpoint
is bound for selective completion and operations are posted without the
completion flag, so the provider increments the counter and writes no queue
entry. In addition, the server's connected endpoints listen on the port after the one
used by the side channel.

\section{Options of the Benchmark}

The benchmark of \cref{sec:sfbench} is linked against either library. Its
options are listed in \cref{tab:sfbench-opts}.

\begin{tightblock}
\begin{longtable}{L{3.4cm}L{11.0cm}}
  \caption[Options of the benchmark]{Options of the benchmark.}\label{tab:sfbench-opts}\\
  \toprule
  Option & Meaning \\
  \midrule
  \endfirsthead
  \toprule
  Option & Meaning \\
  \midrule
  \endhead
  \code{-r} & role, server or client \\
  \code{-H}, \code{-P} & the peer's address, or the address to bind on the server; and the port (19870 by default) \\
  \code{-m} & mode: \code{sig} for the completion signals, \code{bw} for throughput, \code{hol} for the head-of-line experiment \\
  \code{-p}, \code{--hot-prov}, \code{--cold-prov} & provider for both classes, or for one \\
  \code{-t}, \code{--hot-ep}, \code{--cold-ep} & endpoint type, connected or datagram, for both classes or for one \\
  \code{-d} & restrict to one device \\
  \code{-s} & one endpoint for both classes: the head-of-line baseline \\
  \code{--hot-cntr} & hot completions through a counter \\
  \code{--cold-wait} & blocking wait on the cold completion queue \\
  \code{-M} & signal for \code{sig} mode: \code{cq}, \code{flag}, \code{sentinel}, \code{cntr} or \code{read} \\
  \code{--rma-event} & ask the provider to count remote writes \\
  \code{-S} & list of sizes in bytes (8 to 65,536 by default) \\
  \code{-n}, \code{-w}, \code{-W} & iterations (10,000), warm-up iterations (500), and window of outstanding operations (16) \\
  \code{--hot-size}, \code{--cold-size} & sizes of the two streams in \code{hol} mode (64 bytes and 1 MiB) \\
  \code{--hot-gap-us} & pause between hot messages \\
  \code{--no-cold} & run the hot stream alone \\
  \code{--bw-class} & class used in \code{bw} mode \\
  \code{--arena-mb} & size of the registered arena (64 MiB) \\
  \code{--hot-cpu}, \code{--cold-cpu} & processor core for the hot loop and for the cold thread \\
  \code{--perf} & record processor counters \\
  \code{--timeout-ms} & timeout for waits (60,000) \\
  \code{-j} & write results to a JSON file \\
  \bottomrule
\end{longtable}
\end{tightblock}
